\documentclass[11pt,reqno]{amsart}
\usepackage{amsmath,amssymb,amsthm}
\usepackage[margin=1in]{geometry}
\usepackage{booktabs}
\usepackage[colorlinks=true,linkcolor=blue,citecolor=blue,urlcolor=blue]{hyperref}

\theoremstyle{plain}
\newtheorem{theorem}{Theorem}
\newtheorem{proposition}[theorem]{Proposition}
\newtheorem{lemma}[theorem]{Lemma}
\newtheorem{corollary}[theorem]{Corollary}
\theoremstyle{definition}
\newtheorem{definition}[theorem]{Definition}
\theoremstyle{remark}
\newtheorem{remark}[theorem]{Remark}

\newcommand{\R}{\mathbb{R}}
\newcommand{\C}{\mathbb{C}}
\newcommand{\Z}{\mathbb{Z}}
\newcommand{\supp}{\operatorname{supp}}
\newcommand{\eigmin}{\lambda_{\min}}
\DeclareMathOperator{\Lam}{\Lambda}

% DRAFT assembled 2026-06-09 from the research/proof/ document set (00, 00a, 11, 12, 13).
% Every theorem statement below is backed by a written proof in those documents at the
% granularity stated; numerical claims are reproducible via the cited witness scripts
% and gated by scripts/numerics/witness/certify_chain.py.

\title[RH as a tracking inequality]{The Riemann Hypothesis as a tracking inequality:\\
an exact finite-dimensional reduction along the prime filtration\\
{\normalsize\textsc{draft}}}
\author{James Scully}
\address{}
\email{jamesjscully@gmail.com}
\date{June 2026 (draft)}
\subjclass[2020]{11M26 (Primary); 47A57, 47B35, 93C05 (Secondary)}
\keywords{Riemann hypothesis, Weil explicit formula, Weil positivity, Toeplitz forms,
Schur complement, dissipative systems, KYP lemma}

\begin{document}

\begin{abstract}
We reduce the Riemann Hypothesis to a one-parameter family of finite-dimensional scalar
inequalities of tracking type. For each support level $R$, let $M_R = W_\infty - P_R$ be the
Weil form (archimedean minus Euler-prime part) restricted to the pole-killed space of even
real tests supported in $[-R/2,R/2]$. Extending the support across one ``edge'' $L \to R$, we
prove that the prime-power atoms interact with the past through an exactly rank-bounded
channel: the old--new coupling of every atom is a one-sided convolution of edge profiles
(Theorem~\ref{thm:edge}), the old space acts on an edge only through an $r\times r$
compliance matrix $M_c$ and shell matrix $B$ ($r \le 2$ in our bases), and terminal
positivity of the edge is exactly the tracking inequality
$\tau_r = \lambda_{\max}\!\big[M_c^{1/2} W^{*}\,\mathrm{BULK}^{-1}\,W M_c^{1/2}\big] < 1$,
where $W$ measures the deviation of the accumulated arithmetic transgression from a
geometrically determined center (Theorem~\ref{thm:compliance}). RH holds if and only if
$\tau_r < 1$ uniformly along the filtration (Theorem~\ref{thm:equiv}). We then prove the
coordinate is terminal: the tracking residual is identically a cross-correlation of test
data against the nontrivial zeros (Theorem~\ref{thm:zerosum}), so the inequality is
zero-repulsion in exact finite coordinates, and we record a no-go map: eight structural
shortcut classes (subset Euler factors, per-cell estimates, channel decompositions, smooth
surrogates, slack-bearing bounds, among others) are each refuted by explicit finite
counterexamples. The reduction direction is certified against the classical literature
(Weil, Bombieri, Yoshida, Li), including a detection lemma in the constrained even
pole-killed class via an interpolation device of Yoshida. All numerical claims are
reproducible from the accompanying scripts and were validated against the first $5000$
nontrivial zeros. We do not prove RH; we exhibit an exact, finite-dimensional, fail-closed
coordinate system in which its remaining content is a single scalar inequality per edge.
\end{abstract}

\maketitle

\section{Introduction}

Weil's positivity criterion states that RH is equivalent to the positivity of the explicit-formula
functional on a suitable test class \cite{Weil52,Bombieri00,Yoshida92}. Positivity restricted to
tests of small support is unconditional: Connes--Consani \cite{CC20} proved the archimedean form
alone is positive for supports inside $(\tfrac12, 2)$, and Yoshida \cite{Yoshida92} certified
positivity for $a \le \log 2/2$ by finite computation. The difficulty is the passage to all
supports, where the Euler product must carry the positivity.

This paper studies that passage as a \emph{filtration problem}. Writing tests additively
($u = \log x$), we work on the nested family of spaces
\[
H_R = \big\{\, h \in C_c^\infty(\R) \text{ even, real},\ \supp h \subset [-R/2, R/2],\
\hat h(0) = \hat h(i/2) = 0 \,\big\},
\]
and the quadratic form $M_R = W_\infty - P_R$, where $P_R$ is the finite sum of prime-power
Toeplitz atoms reachable at support $R$ and $W_\infty$ the archimedean form. The pole
constraints make $M_R(h,h) = \sum_\rho \hat h\big((\rho-\tfrac12)/i\big)^2$ exactly
(Section~\ref{sec:reduction}), so positivity of $M_R$ for all $R$ is RH.

Our contribution is the exact structure of one step of the filtration. The results are
elementary in form --- support arithmetic, Schur complements, one completed square --- but
they compose into a coordinate system in which the inductive content of RH is isolated in a
single scalar per step, every approximation having been eliminated.

\subsection{Main results}

Fix an edge $0 < L < R$. Decompose $H_R = H_L \oplus S$ (shell $S$), and write the
compressed form in block shape $\begin{pmatrix} A & C \\ C^* & D\end{pmatrix}$.

\begin{theorem}[Edge-window factorization]\label{thm:edge}
Let $x \in H_L$, $y \in H_R$, and let $B_U$ denote the symmetrized translation atom
$B_U(x,y) = \tfrac12\!\int x(v)\,[y(v-U) + y(v+U)]\,dv$. Then:
\begin{enumerate}
\item[(a)] $B_U(x,y) = 0$ for $U > (L+R)/2$ \emph{(dead upper half-edge)};
\item[(b)] for $L \le U \le (L+R)/2$, $B_U(x,y)$ depends on $x$ only through its restriction
to the terminal window $W^* = [\,L/2-(R-L)/2,\ L/2\,]$;
\item[(c)] exactly,
\[
B_U(x,y) \;=\; (q_x * r_y)\big(w(U)\big), \qquad
q_x(s) = x(L/2 - s),\quad r_y(t) = y(-R/2+t),\quad w(U) = \tfrac{L+R}{2} - U ,
\]
a one-sided convolution of edge profiles; consequently the stacked old--shell coupling of all
atoms in $(L,R]$ has rank at most the number of basis bumps meeting $W^*$.
\end{enumerate}
\end{theorem}

\begin{theorem}[Edge-compliance reduction; exact rank-$r$ form]\label{thm:compliance}
Along the chronological chain of one edge the old block is frozen
($E_A^{(j)} = 0$ exactly), and with $U_r$ spanning the coupling column space
(rank $r$ from Theorem~\ref{thm:edge}(c)),
\[
M_c = U_r^* A^{-1} U_r \in \R^{r\times r}, \qquad B = C_0^* A^{-1} U_r, \qquad
V = \textstyle\sum_k T_k \ \ \text{(accumulated transgression data)},
\]
the terminal shell Schur complement satisfies, exactly,
\[
S_N \;=\; \mathrm{BULK}_r \;-\; \big(V + B M_c^{-1}\big)\, M_c\, \big(V + B M_c^{-1}\big)^{*},
\qquad \mathrm{BULK}_r = S_0 + \textstyle\sum E_D + B M_c^{-1} B^{*} .
\]
In particular $S_N \succ 0$ iff $\mathrm{BULK}_r \succ 0$ and
\[
\tau_r \;:=\; \lambda_{\max}\!\Big[\, M_c^{1/2}\, W^{*}\, \mathrm{BULK}_r^{-1}\, W\, M_c^{1/2} \Big]
\;<\; 1, \qquad W = V + B M_c^{-1}.
\]
\end{theorem}

\begin{theorem}[Tracking equivalence]\label{thm:equiv}
RH holds if and only if $\tau_r < 1$ at every edge of the filtration (equivalently,
$M_R \succeq 0$ on $H_R$ for every $R$). The forward direction is
Theorem~\ref{thm:compliance} plus induction along the filtration; the converse is the
detection lemma (Theorem~\ref{thm:detect}).
\end{theorem}

\begin{theorem}[The coordinate is terminal]\label{thm:zerosum}
In the rank-one normalization ($u$ the compliance direction, $\alpha = u^* A^{-1}u$,
$b = C_0^* A^{-1}u$, $v = \sum_k T_k$), the tracking residual obeys the exact identity
$C_N^* A^{-1} u = \alpha\,(v + b/\alpha)$, and on the pole-killed class the form is the
spectral side of the explicit formula; hence
\[
r \;=\; v + \frac{b}{\alpha} \;=\; \frac{2}{\alpha} \sum_{\gamma > 0}
\hat h_{A^{-1}u}(\gamma)\; \hat h_{\mathrm{shell}}(\gamma) ,
\]
a cross-correlation of test data against the nontrivial zeros. The geometric center
$-b/\alpha$ is the explicit-formula prediction of the arithmetic transgression $v$, and the
tracking inequality is zero-repulsion restated.
\end{theorem}

\begin{theorem}[Detection in the constrained class]\label{thm:detect}
Suppose $\zeta$ has a nontrivial zero off the critical line. Then for all sufficiently large
$T$ there exists $h \in H_T$ with $\sum_\rho \hat h\big((\rho-\tfrac12)/i\big)^2 < 0$.
\end{theorem}

The proof of Theorem~\ref{thm:detect} transcribes Yoshida's interpolation device
\cite[Lemma~1, Prop.~1(2)]{Yoshida92} into the pole-killed class via the free
parametrization $\hat h(z) = z^2(z^2+\tfrac14)\hat G(z)$, with two upgrades: a polynomial
envelope of exponent $m \ge 6$ (absorbing the induced weight $|p(z_\rho)|^2 \sim |\rho|^8$
against the Riemann--von Mangoldt count), and a unimodular phase normalization at the target
zero. The construction interpolates rather than concentrates, so no case distinction on the
configuration of off-line zeros arises. An independent second route through Li's criterion
\cite{Li97,BombieriLagarias99,Bombieri00} is recorded alongside.

\subsection{No-go map}\label{sec:nogo-intro}

The coordinate system above was reached by elimination. Each of the following shortcut
classes is refuted by an explicit finite computation (Section~\ref{sec:nogo}): single-place
and subset-Euler passivity; per-cell (atom-local) repair estimates, killed unconditionally by
bounded prime gaps \cite{Maynard15}; smooth (PNT-surrogate) replacements of the atom measure,
which are permanently indefinite; channel decompositions (shift vs.\ autoconvolution), both
channels being indefinite separately; slack-bearing estimates of any kind, since margins decay
like $e^{-1.5R}$ below the natural scale of every term; KMS-weight deformations
$n^{-\beta}$ with $\beta \ne \tfrac12$, which fail on both sides; and frozen-response
closed-form storage in the dissipative-system picture. The Davenport--Heilbronn function,
which satisfies a functional equation but violates its RH, fails the tracking inequality at
the unique edge where its old block is positive ($\tau = 1.554$), and its window form is
indefinite at every later support level --- the criterion separates $\zeta$ from DH.

\subsection*{What this paper does not do}
We do not prove the tracking inequality, and Theorem~\ref{thm:zerosum} is precisely the
reason no proof of it can avoid the zeros: the residual \emph{is} a zero cross-correlation.
The value of the coordinate system is that it is exact (no approximation hypothesis
survives), finite-dimensional per step, certified in both directions against the classical
literature, and instrumented: every claim above is reproducible by the accompanying witness
scripts, and a certification harness re-derives the load-bearing numerics with explicit
pass/fail gates.

\section{Conventions and the reduction}\label{sec:reduction}
% Transcribe from research/proof/00-reduction-to-rh.md:
% - additive coordinates, hat-conventions, the pole-killed space H_R;
% - O1: term-by-term match with Bombieri's printed explicit formula (dictionary
%   f(x) = F(log x) x^{-1/2}); O2 trivial; numerical pin 4.4e-13 / 2000 zeros.
\emph{[To transcribe from the research notes: document 00, sections 2--3.]}

\section{Proof of Theorem~\ref{thm:edge}}

Let $U \ge L$. Consider first $I(U) = \int x(v)\,y(v-U)\,dv$. The integrand is supported
where $v \le L/2$ (support of $x$) and $v - U \ge -R/2$ (support of $y$), i.e.\ on
$[\,U - R/2,\ L/2\,]$, which is nonempty iff $U \le (L+R)/2$. This proves (a) for $I$; the
reflected term $\int x(v)\,y(v+U)\,dv$ is supported on $[-L/2,\ R/2 - U]$, nonempty under the
same condition, and equals $I(U)$ by evenness of $x$ and $y$. The integration window is
contained in $W^* = [\,L/2 - (R-L)/2,\ L/2\,]$ for every such $U$, which proves (b).

For (c), substitute $v = L/2 - s$ and write $U = (L+R)/2 - w$ with $w = w(U) \in [0,(R-L)/2]$:
\[
y(v - U) \;=\; y\!\big(L/2 - s - U\big) \;=\; y\!\big(-R/2 + (w - s)\big) \;=\; r_y(w - s),
\]
and the support condition becomes $0 \le s \le w$. Hence
$I(U) = \int_0^{w} q_x(s)\, r_y(w-s)\, ds = (q_x * r_y)(w)$, exactly.

For the rank bound, fix a local basis $\{x_i\}$ of $H_L$ (bumps centered at $c_i$ with
half-width $a$). The edge profile $q_{x_i}$ vanishes identically on $[0, d_i]$ with
$d_i = L/2 - (c_i + a)$; if $d_i \ge (R-L)/2$ then $(q_{x_i} * r_y)(w) = 0$ for every
admissible $w$, so the corresponding row of the coupling block vanishes for \emph{every} atom
of the edge. The stacked old--shell coupling matrix of all atoms therefore has rank at most
$\#\{\,i : d_i < (R-L)/2\,\}$, the number of bumps meeting $W^*$. \qed

\section{Proof of Theorem~\ref{thm:compliance}}

\emph{Frozen old block.} Every atom of the edge has $U \in (L, R]$; by Theorem~\ref{thm:edge}(a)
applied with $R$ replaced by $L$ (both arguments in $H_L$), $B_U$ vanishes on
$H_L \times H_L$ for $U > L$, and the smooth surrogate restricted to $[L,R]$ vanishes there
for the same support reason. Hence the cell increments satisfy $E_A^{(j)} = 0$ exactly: the
old block $A$ is constant along the chain, and the terminal Schur complement is explicitly
\[
S_N \;=\; D_N - C_N^* A^{-1} C_N, \qquad
D_N = D_0 + \textstyle\sum_k E_D^{(k)}, \qquad
C_N = C_0 + \textstyle\sum_k E_C^{(k)} .
\]

\emph{Rank-$r$ factorization.} By Theorem~\ref{thm:edge}(c) the column space of every
$E_C^{(k)}$ lies in the fixed $r$-dimensional space spanned by the coupled edge-profile
functionals; choose $U_r$ with orthonormal columns spanning it and write
$E_C^{(k)} = U_r T_k^*$. Then $C_N = C_0 + U_r V^*$ with $V = \sum_k T_k$, and expanding,
\[
S_N \;=\; S_0 + \textstyle\sum E_D \;-\; V B^* - B V^* - V M_c V^*,
\qquad S_0 = D_0 - C_0^* A^{-1} C_0 ,
\]
with $M_c = U_r^* A^{-1} U_r$ and $B = C_0^* A^{-1} U_r$. Since $A \succ 0$ (inductive
hypothesis of the filtration) and $U_r$ has full column rank, $M_c \succ 0$. Completing the
square in $V$,
\[
- V B^* - B V^* - V M_c V^* \;=\; B M_c^{-1} B^* \;-\; \big(V + B M_c^{-1}\big) M_c
\big(V + B M_c^{-1}\big)^* ,
\]
which gives the displayed identity with $\mathrm{BULK}_r = S_0 + \sum E_D + B M_c^{-1} B^*$.

\emph{Scalar criterion.} Write $W = V + B M_c^{-1}$ and $Y = W M_c^{1/2}$, so
$S_N = \mathrm{BULK}_r - Y Y^*$. If $S_N \succ 0$ then $\mathrm{BULK}_r \succ Y Y^*
\succeq 0$. Conversely, for $\mathrm{BULK}_r \succ 0$,
\[
S_N \;=\; \mathrm{BULK}_r^{1/2}\Big( I - \mathrm{BULK}_r^{-1/2} Y Y^* \mathrm{BULK}_r^{-1/2}
\Big)\mathrm{BULK}_r^{1/2} \;\succ\; 0
\;\iff\;
\big\| \mathrm{BULK}_r^{-1/2} Y \big\|^2 < 1 ,
\]
and the nonzero eigenvalues of $\mathrm{BULK}_r^{-1/2} Y Y^* \mathrm{BULK}_r^{-1/2}$
coincide with those of the $r \times r$ Gram matrix
$Y^* \mathrm{BULK}_r^{-1} Y = M_c^{1/2} W^* \mathrm{BULK}_r^{-1} W M_c^{1/2}$, whose largest
eigenvalue is $\tau_r$. \qed

\begin{remark}
Numerically, the rank-$1$ and rank-$2$ predictions of $\eigmin(S_N)$ agree with the exact
value to relative error $10^{-7}$--$2\cdot 10^{-5}$ on every measured edge, which is the
Fourier-grid noise floor of the discretization; $\tau_2 < 1$ on every edge, with maximum
$0.954$. See Section~\ref{sec:numerics}.
\end{remark}

\section{Proof of Theorems~\ref{thm:zerosum} and \ref{thm:equiv}}

\emph{The identity.} In the rank-one normalization, $C_N = C_0 + u\,v^*$, hence
\[
C_N^* A^{-1} u \;=\; C_0^* A^{-1} u + v\,(u^* A^{-1} u) \;=\; b + \alpha v
\;=\; \alpha\Big( v + \frac{b}{\alpha} \Big) \;=\; \alpha\, r :
\]
the tracking residual is the terminal old--shell coupling, whitened through $A$, in the
compliance direction.

\emph{The form is the spectral side.} On the pole-killed space the explicit formula
(Section~\ref{sec:reduction}) gives, for $h, h' \in H_R$,
\[
M_R(h, h') \;=\; \sum_\rho \hat h\big(z_\rho\big)\, \hat h'\big(z_\rho\big),
\qquad z_\rho = (\rho - \tfrac12)/i ,
\]
the pole terms vanishing by the constraints and the prime sum truncating at $\log n < R$ by
support. The terminal coupling block $C_N$ is the old--shell compression of $M_R$, so
\[
r \;=\; \frac{1}{\alpha}\, C_N^* A^{-1} u
\;=\; \frac{1}{\alpha} \sum_\rho \hat h_{A^{-1}u}(z_\rho)\, \hat h_{\mathrm{shell}}(z_\rho)
\;=\; \frac{2}{\alpha} \sum_{\gamma > 0} \hat h_{A^{-1}u}(\gamma)\,
\hat h_{\mathrm{shell}}(\gamma)
\]
on RH (zeros pairing as $\pm\gamma$); off RH the same identity holds with the complex
$z_\rho$. This is Theorem~\ref{thm:zerosum}. Numerically the identity closes against the
first $5000$ zeros to relative error $7.7\cdot 10^{-4}$--$3.7\cdot 10^{-3}$ per edge
(Section~\ref{sec:numerics}).

\emph{The equivalence.} If RH holds, $M_R(h,h) = \sum_\rho \hat h(z_\rho)^2 \ge 0$ with real
$z_\rho$, so every compression is positive semidefinite; on the (generic, observed)
positive-definite filtration this yields $\mathrm{BULK}_r \succ 0$ and $\tau_r < 1$ at every
edge by Theorem~\ref{thm:compliance}. Conversely, $\tau_r < 1$ at every edge gives
$S_N \succ 0$ inductively, hence $M_R \succeq 0$ on $H_R$ for all $R$, and
Theorem~\ref{thm:detect} (detection) then forces every zero onto the critical line. This is
Theorem~\ref{thm:equiv}. \qed

\section{Proof of Theorem~\ref{thm:detect}}
% Transcribe research/proof/00a-detection-lemma.md sections 2 (p(D) device), 6''
% (Yoshida interpolation, upgrades U1-U2, symmetrization) and 6''.1 (SL1-SL5).
\emph{[To transcribe from the research notes: document 00a.]}

\section{The no-go map}\label{sec:nogo}
% Transcribe the kill table: research/proof/09 (findings 1-3, refutations),
% 10 (guards G1-G6, S2 channel result), 11 (mass-balance kill), 12 (race data),
% with one witness script per row.
\emph{[To transcribe: the kill table with witness references.]}

\section{Numerical certification}\label{sec:numerics}

All computations use a raised-cosine bump basis ($38$ centers, half-width $0.06$), digamma
kernels at $80$-bit precision, and zeros from \texttt{acb\_zeta\_zero}; scripts are listed at
the section end. A harness re-derives every load-bearing claim from scratch with explicit
pass/fail gates and a nonzero exit code on any flip.

\subsection*{Tracking trajectory ($\zeta$, late edges)}
With $t = -b/\alpha$ the geometric center and $r = v + b/\alpha$ the residual:

\begin{center}
\begin{tabular}{lccccc}
\toprule
edge & $\alpha$ & $\cos(v, t)$ & $|r|/|t|$ & $\tau$ & $\eigmin(S_N)$ \\
\midrule
$4.00 \to 4.26$ & $93.3$ & $-0.03$ & $1.09$ & $0.732$ & $+2.2\cdot 10^{-3}$ \\
$4.26 \to 4.50$ & $168.1$ & $+0.9962$ & $0.181$ & $0.565$ & $+1.6\cdot 10^{-3}$ \\
$4.50 \to 4.76$ & $318.6$ & $+0.9994$ & $0.037$ & $0.179$ & $+2.0\cdot 10^{-3}$ \\
$4.76 \to 4.96$ & $401.4$ & $+0.9995$ & $0.033$ & $0.276$ & $+1.5\cdot 10^{-3}$ \\
\bottomrule
\end{tabular}
\end{center}

The transgression visibly converges to the center, in direction and magnitude, as the
compliance grows. The Davenport--Heilbronn control: at its unique edge with a positive old
block ($1.60 \to 1.88$), $\tau = 1.554 > 1$ and $\eigmin(S_N) = -0.282$; at every later
support level its window form is indefinite ($\eigmin \approx -1.6$). The criterion
discriminates.

\subsection*{The residual is a zero cross-correlation}
Comparing $r$ from cell data against the zero sum of Theorem~\ref{thm:zerosum}
($5000$ zeros, $\gamma_{\max} \approx 5448$):

\begin{center}
\begin{tabular}{lccc}
\toprule
edge & $\mathrm{rel}(r_{\mathrm{chain}}, r_{\mathrm{zeros}})$ &
$\cos(r_{\mathrm{chain}}, r_{\mathrm{zeros}})$ &
$\mathrm{rel}(r_{\mathrm{chain}}, r_{\mathrm{coupling}})$ \\
\midrule
$4.00 \to 4.26$ & $7.7\cdot 10^{-4}$ & $+1.00000$ & $4.7\cdot 10^{-6}$ \\
$4.26 \to 4.50$ & $3.7\cdot 10^{-3}$ & $+1.00000$ & $1.5\cdot 10^{-4}$ \\
$4.50 \to 4.76$ & $2.8\cdot 10^{-3}$ & $+1.00000$ & $1.5\cdot 10^{-4}$ \\
$4.76 \to 4.96$ & $1.2\cdot 10^{-3}$ & $+1.00000$ & $8.8\cdot 10^{-5}$ \\
\bottomrule
\end{tabular}
\end{center}

All discrepancies sit at the zero-truncation/rank-defect scale; the full-window form matches
its zero-sum to relative Frobenius error $8.8\cdot 10^{-6}$.

\subsection*{Harness gates}
\texttt{certify\_chain.py} (all PASS at the time of writing): C1 ladder positivity
($\min \eigmin = +2.9\cdot 10^{-4}$); C2 DH dead from support $1.88$
($\max \eigmin = -7.3\cdot 10^{-2}$); C3 rank-one compliance reproduction on the claimed
range (max.\ rel.\ $2.3\cdot 10^{-5}$); C4 $\tau < 1$ on all living $\zeta$ edges
(max $0.900$ over the full ladder) and DH tracking failure; C5 the zero-sum identity
(max.\ rel.\ $3.7\cdot 10^{-3}$, min.\ cos $1.00000$ against $3000$ zeros).

\subsection*{Scripts}
\texttt{prime\_ladder\_scaling}, \texttt{prime\_edge\_compliance},
\texttt{prime\_edge\_tracking}, \texttt{prime\_edge\_residual\_zerosum},
\texttt{prime\_edge\_rank2\_exact}, \texttt{certify\_chain} in
\texttt{scripts/numerics/witness/}.

\begin{thebibliography}{10}

\bibitem{Weil52} A.~Weil, \emph{Sur les ``formules explicites'' de la th\'eorie des nombres
premiers}, Meddelanden Fr\aa n Lunds Univ. Mat. Sem. (d\'edi\'e \`a M.~Riesz) (1952), 252--265.

\bibitem{Yoshida92} H.~Yoshida, \emph{On Hermitian forms attached to zeta functions}, in:
Zeta Functions in Geometry, Adv. Stud. Pure Math. \textbf{21}, Math. Soc. Japan, Tokyo (1992),
281--325.

\bibitem{Li97} X.-J.~Li, \emph{The positivity of a sequence of numbers and the Riemann
Hypothesis}, J. Number Theory \textbf{65} (1997), 325--333.

\bibitem{BombieriLagarias99} E.~Bombieri and J.~C.~Lagarias, \emph{Complements to Li's
criterion for the Riemann Hypothesis}, J. Number Theory \textbf{77} (1999), 274--287.

\bibitem{Bombieri00} E.~Bombieri, \emph{Remarks on Weil's quadratic functional in the theory
of prime numbers, I}, Rend. Mat. Acc. Lincei (9) \textbf{11} (2000), 183--233.

\bibitem{CC20} A.~Connes and C.~Consani, \emph{Weil positivity and trace formula: the
archimedean place}, Selecta Math. (N.S.) \textbf{27} (2021), no.~77; arXiv:2006.13771.

\bibitem{Maynard15} J.~Maynard, \emph{Small gaps between primes}, Ann. of Math. (2)
\textbf{181} (2015), 383--413.

\bibitem{IK04} H.~Iwaniec and E.~Kowalski, \emph{Analytic Number Theory}, AMS Colloquium
Publications \textbf{53}, 2004.

\end{thebibliography}

\end{document}
